% !TEX root = ../../main.tex
\section{Correlation versus causality}

% (a) Plot births against stork counts.
% (b) Calculate and interpret the first correlation.
% (c) Calculate the area--stork correlation and compare.


\subsection{Birthrate vs. Storks in a Country}

The Datapoints for Storks over the Birthrate is displayed in \autoref{fig:storks-vs-birth}

\begin{figure}[h]
    \centering
    \includegraphics[width=0.4\textwidth]{StorksOverBirth.pdf}
    \caption{Child Birthrate over Storks in a Country}\label{fig:storks-vs-birth}
\end{figure}

\subsection{Correlation Coefficient, Storks vs. Childbirth}\label{sec:s-c}

The Correlation of the two measurement sets has been calculated in MATLAB with \lstinline{corrcoef} and yields
\[
\rho_{S,B} = 0.9162
\]

The high correlation could be explained, because a larger Area of a country has more living room for both humans and storks. More humans would lead to a higher birthrate, so the both variables indirectly linked to each other, so the correlation is spurious.

\subsection{Correlation Coefficient, Storks vs. Country Area}

The correlation coefficient between Storks and Country Area is also positive and quite high with
\[
\rho_{S,A} = 0.5773
\]
which could confirm the observation from \autoref{sec:s-c}.